\section{Introducción}
\label{sec:introduccion}

\subsection{Contexto y Motivación}

Los Large Language Models (LLM) han transformado la inteligencia artificial moderna. Sin embargo, su desarrollo permanece concentrado en pocas organizaciones debido a costes prohibitivos: el entrenamiento de modelos de miles de millones de parámetros requiere infraestructura masiva y decenas de millones de dólares.

Los Small Language Models (SLM) emergen como alternativa viable para organizaciones que requieren modelos controlables, eficientes y especializados. A diferencia de depender de APIs de terceros o servicios en la nube, desarrollar un SLM internamente permite:

\begin{itemize}
    \item Control total sobre datos y privacidad
    \item Optimización para hardware disponible (GPUs consumer-grade)
    \item Especialización en dominios específicos
    \item Reducción de latencia y costes de inferencia
\end{itemize}

El desafío reside en diseñar arquitecturas que capturen las capacidades de modelos grandes manteniendo eficiencia computacional.

\subsection{Problema y Propuesta}

Los Transformers puros, aunque efectivos, requieren múltiples capas de atención para capturar patrones secuenciales, incrementando VRAM y latencia. Las arquitecturas recurrentes (RNN, GRU) son eficientes pero pierden contexto a largo plazo.

Este trabajo propone una \textbf{arquitectura híbrida que combina Transformer Multi-Head Attention y capas GRU dentro de cada bloque}, permitiendo:

\begin{itemize}
    \item Captura de dependencias globales vía Multi-Head Attention
    \item Refinamiento secuencial local vía GRU
    \item Reducción de requisitos de VRAM respecto a Transformers puros equivalentes
\end{itemize}

\subsection{Objetivos}

El objetivo principal es diseñar, implementar y validar esta arquitectura híbrida SLM, evaluando su rendimiento mediante métricas cuantitativas estándar y analizando las limitaciones prácticas del entrenamiento desde cero en entornos con presupuesto de cómputo reducido.

\textbf{Objetivo 1 --- Entrenamiento y Validación Arquitectónica:}
\begin{itemize}
    \item Implementación de la arquitectura híbrida con optimizaciones CUDA
    \item Entrenamiento desde cero sobre datasets conversacionales públicos (OpenAssistant, ShareGPT, Alpaca)
    \item Evaluación cuantitativa mediante métricas estándar de modelado del lenguaje:
    \begin{itemize}
        \item Perplejidad (capacidad predictiva)
        \item Peak VRAM (eficiencia de memoria)
        \item Tokens/segundo (latencia de inferencia)
        \item Precisión Top-1, Top-5 y Top-10
    \end{itemize}
\end{itemize}

\textbf{Objetivo 2 --- Validación de Producción:}
\begin{itemize}
    \item Despliegue en Kubernetes con gestión automatizada mediante ArgoCD
    \item Análisis de cost-benefit: AWS Spot Instances vs desarrollo local
    \item Pipelines de testing automatizados para validación continua
\end{itemize}

\subsection{Contribuciones}

\begin{enumerate}
    \item \textbf{Arquitectura híbrida Transformer-GRU}: diseño e implementación de un bloque híbrido con estados GRU persistentes, validando su convergencia y estabilidad frente a la arquitectura original mediante métricas de gradiente y pérdida
    \item \textbf{Benchmark exhaustivo}: evaluación cuantitativa de 10 métricas sobre 300.000 muestras de validación, con análisis de estabilidad y diagnóstico honesto de las limitaciones cualitativas del modelo resultante
    \item \textbf{Pipeline de ingeniería end-to-end}: implementación completa desde Entrenamiento hasta despliegue en producción, incluyendo infraestructura cloud y automatización de testing
\end{enumerate}

\subsection{Estructura del Documento}

\begin{itemize}
    \item \textbf{Capítulo 2}: Estado del arte --- evolución de LLM a SLM, RNN, LSTM, GRU, Transformers y arquitecturas híbridas modernas
    \item \textbf{Capítulo 3}: Fundamentos teóricos --- formalización matemática del mecanismo de atención, MHA, GRU y análisis de complejidad híbrida
    \item \textbf{Capítulo 4}: Datos y datasets --- descripción de los conjuntos de datos utilizados para el entrenamiento
    \item \textbf{Capítulo 5}: Mejoras arquitectónicas e implementación --- proceso iterativo de desarrollo y decisiones de ingeniería
    \item \textbf{Capítulo 6}: Diseño e implementación de la arquitectura híbrida --- estrategia de mezcla de datos, especificación técnica del modelo y configuración del entrenamiento
    \item \textbf{Capítulo 7}: Evaluación y benchmarking --- 10 métricas de calidad, eficiencia y recursos sobre 300.000 muestras
    \item \textbf{Capítulo 8}: Infraestructura cloud --- ecosistema AWS, hardware G6/G6e/G7e, Terraform, optimización de VRAM y progresión de hardware
    \item \textbf{Capítulo 9}: Análisis de costes --- coste de infraestructura GPU por ciclo de entrenamiento y coste humano equivalente de un ingeniero junior en Data Science
    \item \textbf{Capítulo 10}: MLOps: despliegue y puesta en producción --- CI/CD, reproducibilidad, Rulesets, testing automatizado, Terraform e infraestructura GitOps con Kubernetes y ArgoCD
    \item \textbf{Capítulo 11}: Conclusiones y trabajo futuro
    \item \textbf{Anexo A}: Guía rápida de comandos
    \item \textbf{Anexo B}: Configuración del backend
\end{itemize}